\textbf{iLQR and MPC.} iLQR/DDP in a receding horizon is a standard tool for nonlinear trajectory optimisation~\cite{tassa2012synthesis}; Gauss--Newton and sequential-quadratic-programming views of the algorithm are given in~\cite{roulet2022iterative,abhijeet2025sequential}, and generic tooling in~\cite{andersson2019casadi}. Control limits can be enforced inside the DDP backward pass~\cite{tassa2014control} or with interior-point variants for general inequality constraints~\cite{pavlov2020interior}; our implementation only clips the controls in the forward rollout. Sampling-based MPC is an alternative~\cite{williams2017information}. The design of the stage cost inside MPC has been studied from a suboptimality viewpoint~\cite{beckenbach2019model}; that work adapts the stage cost by learning under stabilising constraints, whereas our energy cost is a different, hand-designed form of stage-cost shaping, without any guarantee.

\textbf{Energy-based swing-up.} The classic controller drives the pendulum energy to its upright value and switches to a stabiliser~\cite{strm2000swinging}. Its gains are hard to choose; tuning by search is studied in~\cite{lee2019enhancement}, and a reachability analysis of a switched energy/LQR cart-pole controller is in~\cite{jaffar2026reachability}. The pendulum as a benchmark is surveyed in~\cite{boubaker2012inverted}.

\textbf{Gap.} In the (limited) literature we looked at, energy shaping appears as a stand-alone feedback law and MPC with quadratic costs. We did not find a controlled head-to-head of an energy-deviation running cost inside iLQR-MPC against both, with weight sweeps, but the search was a few queries and we do not claim that the comparison is new.
